%% start of file `template.tex'.
%% Copyright 2021-2021 Jisoo Jang (simonjisu@gmail.com).
%
% This work may be distributed and/or modified under the
% conditions of the LaTeX Project Public License version 1.3c,
% available at http://www.latex-project.org/lppl/.


\documentclass[11pt,a4paper,roman]{moderncv}        % possible options include font size ('10pt', '11pt' and '12pt'), paper size ('a4paper', 'letterpaper', 'a5paper', 'legalpaper', 'executivepaper' and 'landscape') and font family ('sans' and 'roman')

% modern themes
\moderncvstyle{banking}                            % style options are 'casual' (default), 'classic', 'oldstyle' and 'banking'
\moderncvcolor{blue}                                % color options 'blue' (default), 'orange', 'green', 'red', 'purple', 'grey' and 'black'
% \renewcommand{\familydefault}{\sfdefault}         % to set the default font; use '\sfdefault' for the default sans serif font, '\rmdefault' for the default roman one, or any tex font name
\nopagenumbers{}                                  % uncomment to suppress automatic page numbering for CVs longer than one page

% character encoding
\usepackage[utf8]{inputenc}
\usepackage{fontawesome}
\usepackage{fontspec}
\usepackage{tabularx}
\usepackage{ragged2e}
\usepackage{kotex}
\setmainhangulfont{NanumBarunGothic.ttf}
\usepackage{enumitem}
\setlist[itemize]{labelwidth=0em, labelsep=0.5em, listparindent=1em, itemindent=!, topsep=0.1em, itemsep=0.2em, parsep=0.2em, leftmargin=*, align=left}

% \usepackage{xeCJK}
% \setCJKsansfont{NanumBarunGothic.ttf}
% if you are not using xelatex ou lualatex, replace by the encoding you are using
%\usepackage{CJKutf8}                              % if you need to use CJK to typeset your resume in Chinese, Japanese or Korean

% adjust the page margins
\usepackage[scale=0.85]{geometry}
\usepackage{multicol}
%\setlength{\hintscolumnwidth}{3cm}                % if you want to change the width of the column with the dates
%\setlength{\makecvtitlenamewidth}{10cm}           % for the 'classic' style, if you want to force the width allocated to your name and avoid line breaks. be careful though, the length ishttps://www.overleaf.com/project/60de65a52f1f0ca118dfa8d8 normally calculated to avoid any overlap with your personal info; use this at your own typographical risks...

\usepackage{import}

% personal data
\name{JiSoo}{Jang}
% \title{Curriculum Vitae}                               % optional, remove / comment the line if not wanted
\address{}{}{}% optional, remove / comment the line if not wanted; the "postcode city" and and "country" arguments can be omitted or provided empty
% \phone[mobile]{010-9830-3450}                   % optional, remove / comment the line if not wanted
% \phone[fixed]{01234 123456}                    % optional, remove / comment the line if not wanted
%\phone[fax]{+3~(456)~789~012}                      % optional, remove / comment the line if not wanted
% \email{simonjisu@gmail.com}                               % optional, remove / comment the line if not wanted
% \homepage{https://simonjisu@github.io}                         % optional, remove / comment the line if not wanted
% \extrainfo{}                 % optional, remove / comment the line if not wanted
%\photo[64pt][0.4pt]{picture}                       % optional, remove / comment the line if not wanted; '64pt' is the height the picture must be resized to, 0.4pt is the thickness of the frame around it (put it to 0pt for no frame) and 'picture' is the name of the picture file
%\quote{Some quote}                                 % optional, remove / comment the line if not wanted

% to show numerical labels in the bibliography (default is to show no labels); only useful if you make citations in your resume
%\makeatletter
%\renewcommand*{\bibliographyitemlabel}{\@biblabel{\arabic{enumiv}}}
%\makeatother
%\renewcommand*{\bibliographyitemlabel}{[\arabic{enumiv}]}% CONSIDER REPLACING THE ABOVE BY THIS

% bibliography with mutiple entries
%\usepackage{multibib}
%\newcites{book,misc}{{Books},{Others}}

% \setlength{\parindent}{1.25em}
% \setlength{\parskip}{2em}

\newcommand*{\customcventry}[7][.25em]{
  \begin{tabular}{@{}l} 
    {\bfseries #4}
  \end{tabular}
  \hfill% move it to the right
  \begin{tabular}{l@{}}
     {\bfseries #5}
  \end{tabular} \\
  \begin{tabular}{@{}l} 
    {\itshape #3}
  \end{tabular}
  \hfill% move it to the right
  \begin{tabular}{l@{}}
     {\itshape #2}
  \end{tabular}
  \ifx&#7&%
  \else{\\%
    \begin{minipage}{\maincolumnwidth}%
      \small#7%
    \end{minipage}}\fi%
  \par\addvspace{#1}}

\newcommand*{\customcvproject}[4][.25em]{
%   \vfill\noindent
  \begin{tabular}{@{}l} 
    {\bfseries #2}
  \end{tabular}
  \hfill% move it to the right
  \begin{tabular}{l@{}}
     {\itshape #3}
  \end{tabular}
  \ifx&#4&%
  \else{\\%
    \begin{minipage}{\maincolumnwidth}%
      \small#4%
    \end{minipage}}\fi%
  \par\addvspace{#1}}

\setlength{\tabcolsep}{12pt}

%----------------------------------------------------------------------------------
%            content
%----------------------------------------------------------------------------------
\begin{document}
%\begin{CJK*}{UTF8}{gbsn}                          % to typeset your resume in Chinese using CJK
%-----       resume       ---------------------------------------------------------
\makecvtitle
\vspace*{-23mm}

\begin{center}
\begin{tabular}{ c c c c }
 \faGlobe\enspace \href{https://simonjisu.github.io/}{simonjisu.github.io}  & \faEnvelopeO\enspace simonjisu@snu.ac.kr & \faGithub\enspace \href{https://github.com/simonjisu}{simonjisu} & \faMobile\enspace +82 10-9830-3450\\  
\end{tabular}
\end{center}

\section{RESEARCH}

My research interests lie in the innovative integration of data-driven and model-driven approaches to developing practical applications. I aim to create solutions that are both robust and adaptable. My most recent research involves: 
\begin{itemize}
    \item Designing Agentic AI system architecture for Universal Access to Information
    \item Combining database management with Large Language Models
\end{itemize}

\vspace{5pt}
{\customcvproject{Papers}{}
  {\begin{itemize}
    \item AU-RAG: Agent-based Universal Retrieval Augmented Generation, SIGIR-AP 2024
    \item Prompt Tuning for NL-to-SQL with Embedding Fine-Tuning and RAG, PAKDD 2024 Robust Machine Learning in Open Environments Workshop
    \item Toward Interpretable Machine Learning: Constructing Polynomial Models Based on Feature Interaction Trees, PAKDD 2023
  \end{itemize}
  }
}

\vspace{5pt}
{\customcvproject{Papers(Under Review)}{}
  {\begin{itemize}
    \item Retrieval-Augmented NL2SQL Generation with Data-Centric Query Capsules for Enterprises, SIGIR-AP 2025
    \item Beyond Retrieval: On Demand Derivation and Integration of ML Functions for NL2SQL, ICDE 2026
    \item On Designing and Utilizing a Strategy Map as Navigation Guide for Complex Insights-to-Actions Tasks, AAAI 2026
    \item Behavior Analysis and Context Augmentations for Complex Decision-Making Tasks, AAAI 2026
    \item (Co-Author) Librarian of Tools: Discovery of Compositional Tools for LLM Agents through Embedding-Guided Beam Search, AAAI 2025
  \end{itemize}
  }
}

\vspace{5pt}
{\customcvproject{Awards}{}
  {\begin{itemize}
    \item (24.10) SIGIR Student Travel Grant Award
    \item (23.08) The 11th Ministry of Trade, Industry, and Energy Public Data Utilization BI Contest - 3rd prize in the ``Public Domain'' Area    
  \end{itemize}
  }
}

\section{PROJECTS}
{\customcvproject{Financial Q\&A Chatbot Based on Financial Reports}{2024.11 - 2025.01 (3M)}
    {\begin{itemize}
      \item Building a financial Q\&A chatbot based on financial reports in collaboration between SNU KDT and Naver Cloud
      \item Led a team by helping define the problem, design experiments, and manage the project through weekly meetings
      \item Explained concepts of LLM and RAG, and guided system architecture design by applying the \href{https://arxiv.org/pdf/2312.04511}{LLM Compiler} paper
    \end{itemize}
    }
}

\vspace{5pt}
{\customcvproject{Metric for Deep Learning Visualization Technique}{2021.04 - 2021.06 (2M)}
    {\begin{itemize}
      \item Proposed I-metric, a metric to evaluate neuron visualization as part of Explainable AI techniques
      \item Conducted experiments to verify whether an additional tool (grid-spec activation camera) improves evaluation
      \item Achieved a 28\% improvement over baseline visualization when using the additional tool — \href{http://viplab.snu.ac.kr/viplab/courses/mlvu_2021_1/projects/team18.pdf}{paper link}
      \item Role: Implemented the grid-spec activation camera
    \end{itemize}
    }
}

\vspace{5pt}
{\customcvproject{Improving Cleanness of Dishwasher}{2020.11 - 2021.11 (12M)}
    {\begin{itemize}
      \item A project with SK Networks to improve dishwasher cleaning performance using pattern clustering on turbidity sensor data
    \end{itemize}
    }
}

\vspace{5pt}
{\customcvproject{Battery Life Expectation}{2020.10 - 2021.02 (5M)}
    {\begin{itemize}
      \item Predicted the remaining battery life of a logistics container tracking device by solving it as a regression problem for voltage values
    \end{itemize}
    }
}


\vspace{5pt}
\section{OTHERS}
\begin{minipage}{\maincolumnwidth}%
\small{
\begin{itemize}
  \item \textbf{TA}: 
  \begin{itemize}
    \item (23 Fall - 25 Spring) Prompt Engineering: Art and Science for Interactions with LLM 
    \item (22 Fall - 25 Spring) Project for Data Science
    \item (22 Fall - 23 Spring) Big Data and Knowledge Management Systems 2
  \end{itemize}
  \item \textbf{Lectures}: 
    \begin{itemize}
      \item (19.12.19 - 19.12.21) Introduction to DeepLearning with PyTorch @ Multicampus
      \item (21.07.06 - 21.07.09) Python \& PyTorch Basic @ HanKyong University
      \item (22.07.26 - 22.07.28) Introduction to DeepLearning with PyTorch @ University of Seoul
      \item (24.08.09 - 22.08.09) A Guide to Building LLM-Based Applications: Prompt Engineering and RAG @ Jeonju IT \& CT Industry Promotion Agency
      % LLM 기반 어플리케이션 구축을 위한 가이드: 프롬프트 엔지니어링과 RAG @전주정보문화산업진흥원
    \end{itemize}
  \item \textbf{Conference}: Student Helper to manage the in-person part of the PAKDD 2024 Robust Machine Learning in Open Environments Workshop
  \item \textbf{Publication}: \href{https://product.kyobobook.co.kr/detail/S000001934223}{PyTorch for Deep Learning}
  % 딥러닝에 목마른 사람들을 위한 PyTorch
  
\end{itemize}}%
\end{minipage}%


\vspace{5pt}
\section{EDUCATION}
{\customcventry{2022.03 - current}{Ph.D. Student}{Graduate School of Data Science, Seoul University}{Seoul}{}{}}
{\customcventry{2020.03 - 2022.02}{Master in Data Science}{Graduate School of Data Science, Seoul University}{Seoul}{}{}}
{\customcventry{2010.03 - 2018.08}{BA in Economics}{HanYang University}{Seoul}{}{}}

\vspace{5pt}
\section{EXPERIENCE}
{\customcventry{2018.02 - 2019.08(1.5 Y)}{Boostcourse AI}{Naver Connect}{Seoul}{}
{\begin{itemize}
  \item Publishing AI-related content on \href{https://www.edwith.org/}{edwith} to make it easier for learners to study
  \begin{itemize}
    \item Professor KyungHyun Cho's course: \href{https://www.edwith.org/ai331}{``Natural Language Processing with Deep Learning''}
    \item Professor Andrew Ng's ``Deep Learning'' series: \href{https://www.boostcourse.org/ai215}{1}, \href{https://www.boostcourse.org/ai216}{2}, \href{https://www.boostcourse.org/ai217}{3}, \href{https://www.boostcourse.org/ai218}{4}
  \end{itemize}
  \item Planning and managing project-based learning content
    \begin{itemize}
    \item Assisted in planning and managing the course ``Deep Learning for Everyone Season 2''
    \item Planned, set up the learning environment, and reviewed code for ``Introduction to Deep Learning with \href{https://www.boostcourse.org/ai212}{TensorFlow} and \href{https://www.boostcourse.org/ai214}{PyTorch}''
    \end{itemize}
\end{itemize}

}}

\vspace{5pt}
\section{INTERESTS}
{\customcvproject{Attendances for Learning New Knowledge}{}
  {
  \begin{itemize}
    \item (24.12.09 - 12) Attend the 2nd International ACM SIGIR Conference on Information Retrieval in the Asia Pacific, Tokyo, Japan
    \item (24.05.07 - 10) Attend the 28th Pacific-Asia Conference on Knowledge Discovery and Data Mining, Taipei, Taiwan
    \item (23.05.25 - 28) Attend the 27th Pacific-Asia Conference on Knowledge Discovery and Data Mining, Osaka, Japan
    \item (22.12.17 - 20) Attend 2022 IEEE International Conference on Big Data, Tokyo, Japan  
    \item (18.08.18 - 19) Attend Pycon-Korea 2018, Seoul, Korea
    \item (17.08.12 - 13) Attend Pycon-Korea 2017, Seoul, Korea
    \item (18.01 - 18.04) Participate ``Deep Learning Basic with Pytorch'' Class @ FastCampus, Seoul, Korea
    \item (17.10 - 17.12) Participate ``Text Mining using Python'' Class @ FastCampus, Seoul, Korea
    \item (17.05 - 17.08) Learned ML @ FastCampus Data Science School, Seoul, Korea
  \end{itemize}
  }
}

{\customcvproject{Participate or Running A Study Group}{}
  {\begin{itemize}
    \item (22.11 - 23.01) Head first Rust! @ Modulabs: 
    \item (20.05 - 20.08) Participate Beyond BERT @ Modulabs: \href{https://github.com/modulabs/beyondBERT/issues/20}{Fast-BERT} review
    \item (18.08 - 19.04) Ran a paper study with Yonsei University Master course students (Learning NLP) - \href{https://drive.google.com/drive/folders/1AthPTnvabMrdvFV4Du_lRp3toZ2tk7Pp?usp=sharing}{Drive Link}
  \end{itemize}
  }
}
   
\end{document}
